\documentclass[11pt]{article}

% Packages
\usepackage{amsmath}
\usepackage{amssymb}
\usepackage{amsthm}
\usepackage{geometry}
\usepackage[colorlinks=true, linkcolor=blue, bookmarks=true]{hyperref}
\usepackage{mathrsfs}
\usepackage{mathtools}
\usepackage{tikz}

% Page Geometry
\geometry{letterpaper, margin=1in}

% Theorem Environments
\newtheorem{theorem}{Theorem}
\newtheorem{lemma}[theorem]{Lemma}
\newtheorem{problem}{Problem}
\renewcommand{\qedsymbol}{$\blacksquare$}

% Title Information
\title{\textbf{Advanced Mathematical Logic, Proof Techniques, and Probability Modeling}}
\author{Tashfeen Omran}
\date{September 29, 2026}

\begin{document}

\maketitle

\tableofcontents
\newpage

% ==========================================
% SECTION 1: CONCEPT CHECKLIST & MAPPING
% ==========================================
\section{Concept Checklist \& Question Mapping}
\label{sec:checklist_start}

Based on the rigorous analysis of the provided text, the following theoretical concepts, logical structures, and proof techniques have been extracted and mapped to the generated problem set.

\begin{itemize}
    \item \textbf{C1: Implication Testing via Contrapositive:} Utilizing $P \implies Q \equiv \neg Q \implies \neg P$ to bypass algebraic difficulties. (Problems 1, 2, 3, 4, 5)
    \item \textbf{C2: Domain Exhaustion / Proof by Cases:} Partitioning a domain $\Omega$ into mutually exclusive and collectively exhaustive subsets. (Problems 6, 7, 8, 9, 10)
    \item \textbf{C3: De Morgan's Logical Trick:} Converting intersection requirements into unions via negation. (Problems 11, 12)
    \item \textbf{C4: Factoring to Reveal Parity Sequences:} Utilizing consecutive integer properties (e.g., $n(n+1)$) for modular arithmetic. (Problems 13, 14)
    \item \textbf{C5: Biconditional Splitting:} Proving $P \iff Q$ by independently verifying $P \implies Q$ and $Q \implies P$ (often using contrapositive for one direction). (Problems 15, 16)
    \item \textbf{C6: Proof by Contradiction (Reductio ad Absurdum):} Assuming $P \land \neg Q$ to derive a logical impossibility. (Problems 17, 18, 19, 20)
    \item \textbf{C7: Set-Theoretic Isomorphism \& Containment:} Mapping Boolean logic to topological containment ($P \implies Q$ as $P \subseteq Q$). (Problems 21, 22)
    \item \textbf{C8: Abstract Real-World Modeling (Quant Finance \& ML):} Applying partitions and contrapositives to algorithmic state-spaces, credit risk, and decision trees. (Problems 23, 24, 25)
    \item \textbf{C9: Measure Theory \& Probability Bounds:} Utilizing the Law of Total Probability and Radon-Nikodym analogues over partitioned cases. (Problems 26, 27)
    \item \textbf{C10: Advanced Diagnostic Testing (Find the Flaw):} Identifying vacuous logic errors, converse fallacies, and division-by-zero in cases. (Problems 28, 29, 30)
\end{itemize}

\newpage

% ==========================================
% SECTION 2: PROBLEM SETS
% ==========================================
\section{Problem Sets}
\label{sec:problems}

\subsection{Implication Testing via Contrapositive}
\begin{enumerate}
    \item \textbf{Problem 1:} Let $x \in \mathbb{Z}$. Prove that if $x^3 + 5x + 7$ is even, then $x$ is odd.
    \item \textbf{Problem 2:} Let $x, y \in \mathbb{Z}$. Prove that if $3x - 7y$ is even, then $x$ and $y$ share the same parity.
    \item \textbf{Problem 3:} Let $x, y \in \mathbb{Z}$. Prove that if $x^2 y^2$ is even, then $x$ is even or $y$ is even.
    \item \textbf{Problem 4:} Let $A = \{1, 2, 3\}$ and $B = \{3, 4, 5\}$ be subsets of the universal space $S = \{1, 2, 3, 4, 5, 6\}$. Let $n \in S$. Prove that if $n^2 + 3n + 5$ is even, then $n \in A \cup B$.
    \item \textbf{Problem 5:} Let $x, y, z \in \mathbb{Z}$. Prove that if $x+y$ is odd and $y+z$ is odd, then $x+z$ is even.
\end{enumerate}

\subsection{Domain Exhaustion and Proof by Cases}
\begin{enumerate}
    \setcounter{enumi}{5}
    \item \textbf{Problem 6:} Prove that for any integer $n$, the expression $n^3 - n$ is divisible by 3.
    \item \textbf{Problem 7:} Prove that for any integer $n$, the expression $n^4 - n^2$ is divisible by 4.
    \item \textbf{Problem 8:} Let $x, y \in \mathbb{R}$. Prove rigorously by cases that $\max(x, y) + \min(x, y) = x + y$.
    \item \textbf{Problem 9:} Prove that for any integer $n$, the expression $n^2 + 5n + 13$ is an odd integer.
    \item \textbf{Problem 10:} Prove that if an integer $n$ is not divisible by 3, then $n^2 \equiv 1 \pmod 3$.
\end{enumerate}

\subsection{De Morgan's Tricks and Factoring Sequences}
\begin{enumerate}
    \setcounter{enumi}{10}
    \item \textbf{Problem 11:} Let $A, B, C$ be sets. Prove that $A \setminus (B \cup C) = (A \setminus B) \cap (A \setminus C)$.
    \item \textbf{Problem 12:} Let $x, y \in \mathbb{Z}$. Prove that if $x+y$ is even and $xy$ is odd, then $x$ is odd and $y$ is odd.
    \item \textbf{Problem 13:} Prove that the product of any three consecutive integers is divisible by 6.
    \item \textbf{Problem 14:} Let $n \in \mathbb{Z}$. Prove that $n(n+1)(n+2)(n+3)$ is divisible by 24.
\end{enumerate}

\subsection{Biconditional Splitting}
\begin{enumerate}
    \setcounter{enumi}{14}
    \item \textbf{Problem 15:} Let $x \in \mathbb{Z}$. Prove that $x^3$ is even if and only if $x$ is even.
    \item \textbf{Problem 16:} Let $A$ and $B$ be sets. Prove that $A \setminus B = \emptyset$ if and only if $A \subseteq B$.
\end{enumerate}

\subsection{Proof by Contradiction (Reductio ad Absurdum)}
\begin{enumerate}
    \setcounter{enumi}{16}
    \item \textbf{Problem 17:} Prove that $\sqrt{3}$ is irrational.
    \item \textbf{Problem 18:} Prove that there is no smallest positive rational number.
    \item \textbf{Problem 19:} Prove that the sum of a rational number and an irrational number is strictly irrational.
    \item \textbf{Problem 20:} Prove that the equation $a^2 - 4b = 2$ has no integer solutions for $a, b \in \mathbb{Z}$.
\end{enumerate}

\subsection{Set-Theoretic Isomorphism \& Abstract Modeling}
\begin{enumerate}
    \setcounter{enumi}{20}
    \item \textbf{Problem 21 (Set-Theoretic Logic):} Establish the logical equivalence of the implication $P \implies Q$ and topological containment by proving that $A \subseteq B \iff B^c \subseteq A^c$.
    \item \textbf{Problem 22 (Credit Counterparty Risk):} In a Merton structural default model, Firm A surviving implies Firm B survives. Let $\mathbb{I}_A$ and $\mathbb{I}_B$ be indicator functions for the survival of A and B respectively. Prove that $\mathbb{I}_A \le \mathbb{I}_B$ for all states $\omega \in \Omega$.
    \item \textbf{Problem 23 (Machine Learning - Random Forests):} A decision tree partitions a feature space $\mathbb{R}^2$ into four mutually exclusive leaves $Q_1, Q_2, Q_3, Q_4$ corresponding to the four Cartesian quadrants. Prove that if a data point $(x,y)$ satisfies $(x,y) \notin Q_1 \cup Q_2 \cup Q_3$, it must strictly satisfy $x > 0$ and $y < 0$.
\end{enumerate}

\subsection{Measure Theory \& Probability Bounds}
\begin{enumerate}
    \setcounter{enumi}{23}
    \item \textbf{Problem 24 (Probability Trap - Base Rate Fallacy):} Let $A$ and $B$ be events in a probability space $\Omega$. Prove that $\mathbb{P}(B|A) = 1$ implies $\mathbb{P}(A^c|B^c) = 1$, but provide a counterexample showing it does not imply $\mathbb{P}(A|B) = 1$.
    \item \textbf{Problem 25 (Algorithmic Arbitrage):} Let $\Omega$ be partitioned into three volatility regimes $V_1, V_2, V_3$. Let $X$ be a derivative payoff. If the conditional expectations satisfy $\mathbb{E}[X|V_i] > 0$ for all $i \in \{1,2,3\}$, derive the probability measure to prove that the total expected payoff $\mathbb{E}[X] > 0$.
    \item \textbf{Problem 26 (Measure Theory Intersection):} Let $\mathbb{P}$ be a probability measure on $\Omega$. Prove that if $\mathbb{P}(A) = 1$ and $\mathbb{P}(B) = 1$, then $\mathbb{P}(A \cap B) = 1$.
    \item \textbf{Problem 27 (Stochastic Dominance):} Let $X$ and $Y$ be random variables such that $X(\omega) \le Y(\omega)$ for all $\omega \in \Omega$. Prove by contradiction that $\mathbb{E}[X] \le \mathbb{E}[Y]$.
\end{enumerate}

\subsection{Advanced Diagnostic Testing (Find the Flaw)}
\begin{enumerate}
    \setcounter{enumi}{27}
    \item \textbf{Problem 28 (Vacuous Logic Error):} Find the logical flaw in the following proof: \textit{Theorem: If $0 = 1$, then all integers are odd. Proof: Since $0 \neq 1$, the hypothesis is false. Therefore, the implication is false, meaning there exists an integer that is even.}
    \item \textbf{Problem 29 (Division by Zero in Cases):} Find the flaw in the following algebraic proof by cases: \textit{Theorem: $1 = 2$. Proof: Let $a = b$. Then $a^2 = ab \implies a^2 - b^2 = ab - b^2 \implies (a-b)(a+b) = b(a-b)$. Dividing both sides by $(a-b)$ yields $a+b = b$. Since $a=b$, $2b = b \implies 2 = 1$.}
    \item \textbf{Problem 30 (Converse Fallacy):} Find the flaw in the following proof: \textit{Theorem: If $x^2$ is even, then $x$ is even. Proof: Assume $x$ is even. Then $x = 2k$. Squaring both sides yields $x^2 = 4k^2 = 2(2k^2)$. Since $2k^2$ is an integer, $x^2$ is even. The proof is complete.}
\end{enumerate}

\newpage

% ==========================================
% SECTION 3: DETAILED SOLUTIONS
% ==========================================
\section{Detailed Solutions}
\label{sec:solutions}

\begin{proof}[\textbf{Solution to Problem 1}]
We proceed by contrapositive. The contrapositive statement is: ``If $x$ is even, then $x^3 + 5x + 7$ is odd.''
Assume $x$ is even. By definition, there exists an integer $k$ such that $x = 2k$. Substituting this into the expression yields:
\begin{align*}
x^3 + 5x + 7 &= (2k)^3 + 5(2k) + 7 \\
&= 8k^3 + 10k + 7 \\
&= 8k^3 + 10k + 6 + 1 \\
&= 2(4k^3 + 5k + 3) + 1
\end{align*}
Since $k \in \mathbb{Z}$, the term $(4k^3 + 5k + 3)$ is an integer due to the closure of integers under addition and multiplication. Let $m = 4k^3 + 5k + 3$. The expression becomes $2m + 1$, which is the definition of an odd integer. Since the contrapositive is true, the original implication holds.
\end{proof}

\begin{proof}[\textbf{Solution to Problem 2}]
We proceed by contrapositive. The contrapositive statement is: ``If $x$ and $y$ have opposite parity, then $3x - 7y$ is odd.''
We evaluate the two exhaustive cases for opposite parity.
\textbf{Case 1:} $x$ is even and $y$ is odd. Let $x = 2k$ and $y = 2m + 1$ for $k, m \in \mathbb{Z}$.
\begin{align*}
3x - 7y &= 3(2k) - 7(2m + 1) \\
&= 6k - 14m - 7 \\
&= 6k - 14m - 8 + 1 \\
&= 2(3k - 7m - 4) + 1
\end{align*}
Since $(3k - 7m - 4) \in \mathbb{Z}$, the result is odd.
\textbf{Case 2:} $x$ is odd and $y$ is even. Let $x = 2k + 1$ and $y = 2m$ for $k, m \in \mathbb{Z}$.
\begin{align*}
3x - 7y &= 3(2k + 1) - 7(2m) \\
&= 6k + 3 - 14m \\
&= 6k - 14m + 2 + 1 \\
&= 2(3k - 7m + 1) + 1
\end{align*}
Since $(3k - 7m + 1) \in \mathbb{Z}$, the result is odd.
In all cases, the contrapositive holds. Thus, the original implication is true.
\end{proof}

\begin{proof}[\textbf{Solution to Problem 3}]
We proceed by contrapositive. By De Morgan's Laws, the negation of ``$x$ is even or $y$ is even'' is ``$x$ is odd and $y$ is odd.'' The contrapositive is: ``If $x$ is odd and $y$ is odd, then $x^2 y^2$ is odd.''
Assume $x$ and $y$ are odd. Let $x = 2k + 1$ and $y = 2m + 1$ for $k, m \in \mathbb{Z}$.
First, we find $xy$:
\begin{align*}
xy &= (2k + 1)(2m + 1) \\
&= 4km + 2k + 2m + 1 \\
&= 2(2km + k + m) + 1
\end{align*}
Let $p = 2km + k + m$. Since $p \in \mathbb{Z}$, $xy = 2p + 1$, meaning $xy$ is odd.
Now we evaluate $x^2 y^2 = (xy)^2$:
\begin{align*}
(xy)^2 &= (2p + 1)^2 \\
&= 4p^2 + 4p + 1 \\
&= 2(2p^2 + 2p) + 1
\end{align*}
Since $(2p^2 + 2p) \in \mathbb{Z}$, $x^2 y^2$ is odd. The contrapositive is proven.
\end{proof}

\begin{proof}[\textbf{Solution to Problem 4}]
We proceed by contrapositive: ``If $n \notin A \cup B$, then $n^2 + 3n + 5$ is odd.''
First, we determine the set $A \cup B$.
\[ A \cup B = \{1, 2, 3\} \cup \{3, 4, 5\} = \{1, 2, 3, 4, 5\} \]
The domain is $S = \{1, 2, 3, 4, 5, 6\}$. The condition $n \notin A \cup B$ within $S$ leaves exactly one case: $n = 6$.
We evaluate the expression for $n = 6$:
\[ n^2 + 3n + 5 = (6)^2 + 3(6) + 5 = 36 + 18 + 5 = 59 \]
Since $59 = 2(29) + 1$, it is an odd integer. The contrapositive holds, verifying the original implication.
\end{proof}

\begin{proof}[\textbf{Solution to Problem 5}]
We use a direct proof. Assume $x+y$ is odd and $y+z$ is odd. By definition, there exist integers $k, m$ such that:
\[ x + y = 2k + 1 \]
\[ y + z = 2m + 1 \]
We add the two equations together:
\begin{align*}
(x + y) + (y + z) &= (2k + 1) + (2m + 1) \\
x + z + 2y &= 2k + 2m + 2
\end{align*}
We isolate $x + z$:
\begin{align*}
x + z &= 2k + 2m + 2 - 2y \\
x + z &= 2(k + m + 1 - y)
\end{align*}
Since $k, m, y \in \mathbb{Z}$, the term $(k + m + 1 - y)$ is an integer. Therefore, $x+z$ is even.
\end{proof}

\begin{proof}[\textbf{Solution to Problem 6}]
We proceed by cases based on the Division Algorithm. Any integer $n$ can be expressed in one of three forms modulo 3: $n = 3k$, $n = 3k+1$, or $n = 3k+2$ for some $k \in \mathbb{Z}$. We factor the expression: $n^3 - n = n(n^2 - 1) = n(n-1)(n+1)$.
\textbf{Case 1:} $n = 3k$.
\[ n^3 - n = (3k)(3k-1)(3k+1) = 3[k(3k-1)(3k+1)] \]
This is a multiple of 3.
\textbf{Case 2:} $n = 3k + 1$.
\[ n^3 - n = (3k+1)(3k)(3k+2) = 3[k(3k+1)(3k+2)] \]
This is a multiple of 3.
\textbf{Case 3:} $n = 3k + 2$.
\[ n^3 - n = (3k+2)(3k+1)(3k+3) = (3k+2)(3k+1)3(k+1) = 3[(3k+2)(3k+1)(k+1)] \]
This is a multiple of 3.
Since all exhaustive cases yield a multiple of 3, $n^3 - n$ is divisible by 3 for all $n \in \mathbb{Z}$.
\end{proof}

\begin{proof}[\textbf{Solution to Problem 7}]
We proceed by cases on the parity of $n$. We factor the expression: $n^4 - n^2 = n^2(n^2 - 1)$.
\textbf{Case 1:} $n$ is even ($n = 2k$).
\[ n^4 - n^2 = (2k)^2((2k)^2 - 1) = 4k^2(4k^2 - 1) = 4[k^2(4k^2 - 1)] \]
Since the term in brackets is an integer, the expression is divisible by 4.
\textbf{Case 2:} $n$ is odd ($n = 2k + 1$).
\begin{align*}
n^4 - n^2 &= (2k+1)^2((2k+1)^2 - 1) \\
&= (2k+1)^2(4k^2 + 4k + 1 - 1) \\
&= (2k+1)^2(4k^2 + 4k) \\
&= 4[k(k+1)(2k+1)^2]
\end{align*}
Since the term in brackets is an integer, the expression is divisible by 4. Both cases hold.
\end{proof}

\begin{proof}[\textbf{Solution to Problem 8}]
We proceed by cases based on the relative order of $x$ and $y$. By the Axiom of Trichotomy, either $x \ge y$ or $x < y$.
\textbf{Case 1:} $x \ge y$.
By definition, $\max(x, y) = x$ and $\min(x, y) = y$.
\[ \max(x, y) + \min(x, y) = x + y \]
\textbf{Case 2:} $x < y$.
By definition, $\max(x, y) = y$ and $\min(x, y) = x$.
\[ \max(x, y) + \min(x, y) = y + x = x + y \]
Since the equality holds in all collectively exhaustive cases, the theorem is proven.
\end{proof}

\begin{proof}[\textbf{Solution to Problem 9}]
We proceed by cases on the parity of $n$.
\textbf{Case 1:} $n$ is even ($n = 2k$).
\begin{align*}
n^2 + 5n + 13 &= (2k)^2 + 5(2k) + 13 \\
&= 4k^2 + 10k + 12 + 1 \\
&= 2(2k^2 + 5k + 6) + 1
\end{align*}
This is odd.
\textbf{Case 2:} $n$ is odd ($n = 2k + 1$).
\begin{align*}
n^2 + 5n + 13 &= (2k+1)^2 + 5(2k+1) + 13 \\
&= 4k^2 + 4k + 1 + 10k + 5 + 13 \\
&= 4k^2 + 14k + 19 \\
&= 4k^2 + 14k + 18 + 1 \\
&= 2(2k^2 + 7k + 9) + 1
\end{align*}
This is odd. The expression is odd for all $n \in \mathbb{Z}$.
\end{proof}

\begin{proof}[\textbf{Solution to Problem 10}]
If $n$ is not divisible by 3, then by the Division Algorithm, $n$ must leave a remainder of 1 or 2 when divided by 3. We proceed by cases.
\textbf{Case 1:} $n = 3k + 1$ for $k \in \mathbb{Z}$.
\begin{align*}
n^2 &= (3k + 1)^2 \\
&= 9k^2 + 6k + 1 \\
&= 3(3k^2 + 2k) + 1
\end{align*}
Thus, $n^2 \equiv 1 \pmod 3$.
\textbf{Case 2:} $n = 3k + 2$ for $k \in \mathbb{Z}$.
\begin{align*}
n^2 &= (3k + 2)^2 \\
&= 9k^2 + 12k + 4 \\
&= 9k^2 + 12k + 3 + 1 \\
&= 3(3k^2 + 4k + 1) + 1
\end{align*}
Thus, $n^2 \equiv 1 \pmod 3$. Both cases satisfy the theorem.
\end{proof}

\begin{proof}[\textbf{Solution to Problem 11}]
We use set-builder notation and De Morgan's Laws.
\begin{align*}
A \setminus (B \cup C) &= \{ x \mid x \in A \land x \notin (B \cup C) \} \\
&= \{ x \mid x \in A \land \neg(x \in B \lor x \in C) \} \\
&= \{ x \mid x \in A \land (x \notin B \land x \notin C) \} \quad \text{(De Morgan's Law)} \\
&= \{ x \mid (x \in A \land x \notin B) \land (x \in A \land x \notin C) \} \quad \text{(Idempotence of } \land \text{)} \\
&= \{ x \mid x \in (A \setminus B) \land x \in (A \setminus C) \} \\
&= (A \setminus B) \cap (A \setminus C)
\end{align*}
\end{proof}

\begin{proof}[\textbf{Solution to Problem 12}]
We proceed by contrapositive. The negation of ``$x$ is odd and $y$ is odd'' is ``$x$ is even or $y$ is even'' (De Morgan's Law). The contrapositive is: ``If $x$ is even or $y$ is even, then $x+y$ is odd or $xy$ is even.''
Assume $x$ is even or $y$ is even. We proceed by cases.
\textbf{Case 1:} $x$ is even ($x = 2k$). Let $y \in \mathbb{Z}$ be arbitrary.
Then $xy = (2k)y = 2(ky)$. Since $ky \in \mathbb{Z}$, $xy$ is even. Thus, the disjunction ``$x+y$ is odd or $xy$ is even'' is true.
\textbf{Case 2:} $y$ is even ($y = 2m$). Let $x \in \mathbb{Z}$ be arbitrary.
Then $xy = x(2m) = 2(xm)$. Since $xm \in \mathbb{Z}$, $xy$ is even. The disjunction holds.
In all cases, the contrapositive is true.
\end{proof}

\begin{proof}[\textbf{Solution to Problem 13}]
Let the three consecutive integers be $n, n+1, n+2$. Their product is $P = n(n+1)(n+2)$.
To prove $P$ is divisible by 6, we must prove it is divisible by 2 and by 3 (since $\gcd(2,3)=1$).
\textbf{Divisibility by 2:} Among any two consecutive integers $n$ and $n+1$, exactly one must be even. Thus, their product $n(n+1)$ is even, meaning $P$ is even (divisible by 2).
\textbf{Divisibility by 3:} By the Division Algorithm, $n$ takes the form $3k, 3k+1,$ or $3k+2$.
If $n = 3k$, $3 \mid n$, so $3 \mid P$.
If $n = 3k+1$, then $n+2 = 3k+3 = 3(k+1)$, so $3 \mid (n+2)$, meaning $3 \mid P$.
If $n = 3k+2$, then $n+1 = 3k+3 = 3(k+1)$, so $3 \mid (n+1)$, meaning $3 \mid P$.
Since $P$ is divisible by both 2 and 3, it is divisible by 6.
\end{proof}

\begin{proof}[\textbf{Solution to Problem 14}]
Let $P = n(n+1)(n+2)(n+3)$. We must show $24 \mid P$. Since $24 = 8 \times 3$ and $\gcd(8,3)=1$, we show $P$ is divisible by 3 and by 8.
\textbf{Divisibility by 3:} $P$ contains 4 consecutive integers. By the Pigeonhole Principle, at least one must be a multiple of 3. Thus, $3 \mid P$.
\textbf{Divisibility by 8:} Among four consecutive integers, exactly two are even. Let them be $2k$ and $2k+2$. Their product is $(2k)(2k+2) = 4k(k+1)$. Since $k$ and $k+1$ are consecutive, one of them is even, meaning $k(k+1)$ is divisible by 2. Thus, $4k(k+1)$ is divisible by $4 \times 2 = 8$. Therefore, $8 \mid P$.
Since $3 \mid P$ and $8 \mid P$, $24 \mid P$.
\end{proof}

\begin{proof}[\textbf{Solution to Problem 15}]
We must prove both directions of the biconditional.
($\implies$): We prove the contrapositive: ``If $x$ is odd, then $x^3$ is odd.''
Assume $x = 2k + 1$.
\begin{align*}
x^3 &= (2k + 1)^3 \\
&= 8k^3 + 12k^2 + 6k + 1 \\
&= 2(4k^3 + 6k^2 + 3k) + 1
\end{align*}
This is odd.
($\impliedby$): We use a direct proof. Assume $x$ is even, so $x = 2k$.
\[ x^3 = (2k)^3 = 8k^3 = 2(4k^3) \]
This is even. Both directions hold.
\end{proof}

\begin{proof}[\textbf{Solution to Problem 16}]
We prove both directions.
($\implies$): Assume $A \setminus B = \emptyset$. We must show $A \subseteq B$. Let $x \in A$. Suppose, for the sake of contradiction, that $x \notin B$. Then by definition, $x \in A \setminus B$. But $A \setminus B = \emptyset$, meaning it contains no elements. This is a contradiction. Therefore, $x \in B$, establishing $A \subseteq B$.
($\impliedby$): Assume $A \subseteq B$. We must show $A \setminus B = \emptyset$. Suppose, for the sake of contradiction, that $A \setminus B \neq \emptyset$. Then there exists some $x \in A \setminus B$. By definition, $x \in A$ and $x \notin B$. However, since $A \subseteq B$, $x \in A \implies x \in B$. This contradicts $x \notin B$. Thus, $A \setminus B = \emptyset$.
\end{proof}

\begin{proof}[\textbf{Solution to Problem 17}]
Assume, for the sake of contradiction, that $\sqrt{3}$ is rational. Then it can be expressed as a fraction in lowest terms: $\sqrt{3} = \frac{a}{b}$, where $a, b \in \mathbb{Z}$, $b \neq 0$, and $\gcd(a, b) = 1$.
Squaring both sides yields $3 = \frac{a^2}{b^2}$, which implies $a^2 = 3b^2$.
This means $a^2$ is a multiple of 3. By Euclid's Lemma (since 3 is prime), $a$ must also be a multiple of 3. Let $a = 3k$ for some $k \in \mathbb{Z}$.
Substitute $a$ back into the equation:
\[ (3k)^2 = 3b^2 \implies 9k^2 = 3b^2 \implies b^2 = 3k^2 \]
This implies $b^2$ is a multiple of 3, which means $b$ must also be a multiple of 3.
We have deduced that both $a$ and $b$ are multiples of 3, contradicting our assumption that $\gcd(a, b) = 1$. Therefore, $\sqrt{3}$ is irrational.
\end{proof}

\begin{proof}[\textbf{Solution to Problem 18}]
Assume, for the sake of contradiction, that there exists a smallest positive rational number, say $r$.
Since $r$ is rational and positive, $r = \frac{a}{b}$ where $a, b \in \mathbb{Z}^+$.
Consider the number $q = \frac{r}{2}$. Since $r > 0$, $q > 0$.
Furthermore, $q = \frac{a}{2b}$. Since $a, 2b \in \mathbb{Z}^+$, $q$ is a positive rational number.
However, $q = \frac{r}{2} < r$. We have found a positive rational number strictly smaller than $r$, contradicting the assumption that $r$ is the smallest positive rational number. Thus, no such number exists.
\end{proof}

\begin{proof}[\textbf{Solution to Problem 19}]
Let $q$ be a rational number and $x$ be an irrational number. Assume, for the sake of contradiction, that their sum is rational. Let $q + x = p$, where $p \in \mathbb{Q}$.
Since $q \in \mathbb{Q}$ and $p \in \mathbb{Q}$, they can be written as $q = \frac{a}{b}$ and $p = \frac{c}{d}$ for integers $a,b,c,d$ with $b,d \neq 0$.
We isolate $x$:
\[ x = p - q = \frac{c}{d} - \frac{a}{b} = \frac{cb - ad}{db} \]
Since $cb - ad \in \mathbb{Z}$ and $db \in \mathbb{Z}$ (with $db \neq 0$), $x$ is a rational number. This contradicts the premise that $x$ is irrational. Therefore, the sum must be irrational.
\end{proof}

\begin{proof}[\textbf{Solution to Problem 20}]
Assume, for the sake of contradiction, that there exist integers $a, b$ such that $a^2 - 4b = 2$.
We isolate $a^2$:
\[ a^2 = 4b + 2 = 2(2b + 1) \]
This implies $a^2$ is even, which means $a$ must be even. Let $a = 2k$ for some $k \in \mathbb{Z}$.
Substitute $a$ back into the equation:
\begin{align*}
(2k)^2 &= 2(2b + 1) \\
4k^2 &= 4b + 2 \\
2k^2 &= 2b + 1
\end{align*}
The left side of the equation, $2k^2$, is an even integer. The right side, $2b + 1$, is an odd integer. An integer cannot be both even and odd. This is a contradiction. Thus, no such integer solutions exist.
\end{proof}

\begin{proof}[\textbf{Solution to Problem 21}]
We must prove $A \subseteq B \iff B^c \subseteq A^c$.
($\implies$): Assume $A \subseteq B$. Let $x \in B^c$. By definition of the complement, $x \notin B$. Since $A \subseteq B$, if $x$ were in $A$, it would have to be in $B$. Since $x \notin B$, it must be that $x \notin A$. Thus, $x \in A^c$. This proves $B^c \subseteq A^c$.
($\impliedby$): Assume $B^c \subseteq A^c$. Let $x \in A$. Suppose, for contradiction, that $x \notin B$. Then $x \in B^c$. By our assumption, $x \in B^c \implies x \in A^c$, meaning $x \notin A$. This contradicts $x \in A$. Thus, $x \in B$, proving $A \subseteq B$.
\end{proof}

\begin{proof}[\textbf{Solution to Problem 22}]
Let $A$ be the event ``Firm A survives'' and $B$ be the event ``Firm B survives''. The premise ``Firm A surviving implies Firm B survives'' means $A \subseteq B$.
The indicator function $\mathbb{I}_E(\omega)$ equals 1 if $\omega \in E$ and 0 otherwise.
We proceed by cases on $\omega \in \Omega$.
\textbf{Case 1:} $\omega \in A$. Since $A \subseteq B$, $\omega \in B$. Thus, $\mathbb{I}_A(\omega) = 1$ and $\mathbb{I}_B(\omega) = 1$. Here, $1 \le 1$ holds.
\textbf{Case 2:} $\omega \notin A$. Then $\mathbb{I}_A(\omega) = 0$. Since $\mathbb{I}_B(\omega) \in \{0, 1\}$, $0 \le \mathbb{I}_B(\omega)$ trivially holds.
In all cases, $\mathbb{I}_A \le \mathbb{I}_B$.
\end{proof}

\begin{proof}[\textbf{Solution to Problem 23}]
Let the quadrants be defined as:
$Q_1 = \{(x,y) \mid x > 0, y > 0\}$
$Q_2 = \{(x,y) \mid x < 0, y > 0\}$
$Q_3 = \{(x,y) \mid x < 0, y < 0\}$
$Q_4 = \{(x,y) \mid x > 0, y < 0\}$
(Assuming boundaries are handled such that they partition $\mathbb{R}^2$).
The premise is $(x,y) \notin Q_1 \cup Q_2 \cup Q_3$. By De Morgan's Laws for Sets:
\[ (Q_1 \cup Q_2 \cup Q_3)^c = Q_1^c \cap Q_2^c \cap Q_3^c \]
Since $Q_1, Q_2, Q_3, Q_4$ form a partition of $\mathbb{R}^2$, the complement of the union of any three is exactly the fourth. Thus, $(Q_1 \cup Q_2 \cup Q_3)^c = Q_4$.
Therefore, $(x,y) \in Q_4$, which by definition strictly satisfies $x > 0$ and $y < 0$.

\begin{center}
\begin{tikzpicture}[scale=0.8]
    \draw[thick, ->] (-3,0) -- (3,0) node[right] {$x$};
    \draw[thick, ->] (0,-3) -- (0,3) node[above] {$y$};
    \node at (1.5, 1.5) {$Q_1$};
    \node at (-1.5, 1.5) {$Q_2$};
    \node at (-1.5, -1.5) {$Q_3$};
    \node at (1.5, -1.5) {\textbf{$Q_4$}};
\end{tikzpicture}
\end{center}
\end{proof}

\begin{proof}[\textbf{Solution to Problem 24}]
First, we prove $\mathbb{P}(B|A) = 1 \implies \mathbb{P}(A^c|B^c) = 1$.
By definition, $\mathbb{P}(B|A) = \frac{\mathbb{P}(A \cap B)}{\mathbb{P}(A)} = 1 \implies \mathbb{P}(A \cap B) = \mathbb{P}(A)$.
This implies $\mathbb{P}(A \setminus B) = \mathbb{P}(A) - \mathbb{P}(A \cap B) = 0$.
We want to find $\mathbb{P}(A^c|B^c) = \frac{\mathbb{P}(A^c \cap B^c)}{\mathbb{P}(B^c)}$.
Note that $B^c = (B^c \cap A) \cup (B^c \cap A^c)$. Thus, $\mathbb{P}(B^c) = \mathbb{P}(A \setminus B) + \mathbb{P}(A^c \cap B^c) = 0 + \mathbb{P}(A^c \cap B^c)$.
Therefore, $\mathbb{P}(A^c|B^c) = \frac{\mathbb{P}(A^c \cap B^c)}{\mathbb{P}(A^c \cap B^c)} = 1$.

\textbf{Counterexample for $\mathbb{P}(A|B) = 1$:} Let $\Omega = \{1, 2, 3\}$ with uniform probability. Let $A = \{1\}$ and $B = \{1, 2\}$.
$\mathbb{P}(B|A) = \frac{\mathbb{P}(\{1\})}{\mathbb{P}(\{1\})} = 1$.
However, $\mathbb{P}(A|B) = \frac{\mathbb{P}(\{1\})}{\mathbb{P}(\{1, 2\})} = \frac{1/3}{2/3} = 0.5 \neq 1$.
\end{proof}

\begin{proof}[\textbf{Solution to Problem 25}]
By the definition of a partition, $\bigcup_{i=1}^3 V_i = \Omega$ and $V_i \cap V_j = \emptyset$ for $i \neq j$.
By the Law of Total Probability (discrete Radon-Nikodym extension), the total expectation is the sum of conditional expectations weighted by their probabilities:
\[ \mathbb{E}[X] = \sum_{i=1}^3 \mathbb{E}[X | V_i] \mathbb{P}(V_i) \]
We are given $\mathbb{E}[X | V_i] > 0$ for all $i$.
Since $\mathbb{P}$ is a probability measure, $\mathbb{P}(V_i) \ge 0$ for all $i$, and $\sum_{i=1}^3 \mathbb{P}(V_i) = 1$.
Because the sum of the probabilities is 1, at least one partition must have strictly positive probability, say $\mathbb{P}(V_k) > 0$.
Thus, the sum consists of non-negative terms, with at least one strictly positive term:
\[ \mathbb{E}[X | V_k] \mathbb{P}(V_k) > 0 \]
Therefore, the entire sum $\mathbb{E}[X] > 0$.
\end{proof}

\begin{proof}[\textbf{Solution to Problem 26}]
We use the properties of probability measures and De Morgan's Laws.
Given $\mathbb{P}(A) = 1$, we have $\mathbb{P}(A^c) = 1 - 1 = 0$.
Given $\mathbb{P}(B) = 1$, we have $\mathbb{P}(B^c) = 1 - 1 = 0$.
We want to find $\mathbb{P}(A \cap B)$. By the complement rule:
\[ \mathbb{P}(A \cap B) = 1 - \mathbb{P}((A \cap B)^c) \]
By De Morgan's Law, $(A \cap B)^c = A^c \cup B^c$.
By the Union Bound (Boole's Inequality):
\[ \mathbb{P}(A^c \cup B^c) \le \mathbb{P}(A^c) + \mathbb{P}(B^c) = 0 + 0 = 0 \]
Since probability cannot be negative, $\mathbb{P}(A^c \cup B^c) = 0$.
Thus, $\mathbb{P}(A \cap B) = 1 - 0 = 1$.
\end{proof}

\begin{proof}[\textbf{Solution to Problem 27}]
Assume, for the sake of contradiction, that $\mathbb{E}[X] > \mathbb{E}[Y]$.
By linearity of expectation, this implies $\mathbb{E}[X] - \mathbb{E}[Y] > 0 \implies \mathbb{E}[X - Y] > 0$.
However, we are given that $X(\omega) \le Y(\omega)$ for all $\omega \in \Omega$.
This means $X(\omega) - Y(\omega) \le 0$ for all $\omega \in \Omega$.
The expectation of a non-positive random variable must be non-positive. That is, since $X - Y \le 0$ almost surely, $\mathbb{E}[X - Y] \le 0$.
This directly contradicts our derived statement $\mathbb{E}[X - Y] > 0$. Thus, the assumption must be false, and $\mathbb{E}[X] \le \mathbb{E}[Y]$.
\end{proof}

\begin{proof}[\textbf{Solution to Problem 28 (Diagnostic)}]
The flaw is a fundamental misunderstanding of \textbf{vacuous truth}. In formal logic, an implication $P \implies Q$ is defined as $\neg P \lor Q$. If the hypothesis $P$ is false (as $0 = 1$ is false), then $\neg P$ is true, making the entire implication statement logically \textbf{True}, regardless of the conclusion $Q$. The student incorrectly claimed that a false hypothesis makes the implication false. Because the implication is vacuously true, it does not guarantee the physical reality of the conclusion (that all integers are odd), nor does it imply the negation of the conclusion.
\end{proof}

\begin{proof}[\textbf{Solution to Problem 29 (Diagnostic)}]
The flaw occurs in the step: ``Dividing both sides by $(a-b)$ yields $a+b = b$.''
Because the premise establishes $a = b$, it follows that $a - b = 0$. Dividing by $(a-b)$ is equivalent to dividing by zero, which is undefined and violates the axioms of real numbers. A rigorous proof by cases must isolate division by zero. If one sets up Case 1: $a - b \neq 0$, the division is valid, but contradicts the premise $a=b$. Case 2: $a - b = 0$, the equation becomes $0 = 0$, which provides no information to deduce $2=1$.
\end{proof}

\begin{proof}[\textbf{Solution to Problem 30 (Diagnostic)}]
The flaw is the \textbf{Fallacy of the Converse} (Affirming the Consequent). The theorem states $P \implies Q$ (If $x^2$ is even, then $x$ is even). The student assumed $Q$ ($x$ is even) and proved $P$ ($x^2$ is even). Proving $Q \implies P$ does not establish $P \implies Q$. To correctly prove this theorem, the student should have used the contrapositive: Assume $x$ is odd, and prove $x^2$ is odd.
\end{proof}

\newpage

% ==========================================
% SECTION 4: CONCEPT CHECKLIST (REPEATED)
% ==========================================
\section{Concept Checklist \& Question Mapping (Repeated)}
\label{sec:checklist_end}

As requested, here is the repeated reference mapping to ensure all theoretical concepts from the PDF were rigorously tested:

\begin{itemize}
    \item \textbf{C1: Implication Testing via Contrapositive:} Utilizing $P \implies Q \equiv \neg Q \implies \neg P$ to bypass algebraic difficulties. (Problems 1, 2, 3, 4, 5)
    \item \textbf{C2: Domain Exhaustion / Proof by Cases:} Partitioning a domain $\Omega$ into mutually exclusive and collectively exhaustive subsets. (Problems 6, 7, 8, 9, 10)
    \item \textbf{C3: De Morgan's Logical Trick:} Converting intersection requirements into unions via negation. (Problems 11, 12)
    \item \textbf{C4: Factoring to Reveal Parity Sequences:} Utilizing consecutive integer properties (e.g., $n(n+1)$) for modular arithmetic. (Problems 13, 14)
    \item \textbf{C5: Biconditional Splitting:} Proving $P \iff Q$ by independently verifying $P \implies Q$ and $Q \implies P$ (often using contrapositive for one direction). (Problems 15, 16)
    \item \textbf{C6: Proof by Contradiction (Reductio ad Absurdum):} Assuming $P \land \neg Q$ to derive a logical impossibility. (Problems 17, 18, 19, 20)
    \item \textbf{C7: Set-Theoretic Isomorphism \& Containment:} Mapping Boolean logic to topological containment ($P \implies Q$ as $P \subseteq Q$). (Problems 21, 22)
    \item \textbf{C8: Abstract Real-World Modeling (Quant Finance \& ML):} Applying partitions and contrapositives to algorithmic state-spaces, credit risk, and decision trees. (Problems 23, 24, 25)
    \item \textbf{C9: Measure Theory \& Probability Bounds:} Utilizing the Law of Total Probability and Radon-Nikodym analogues over partitioned cases. (Problems 26, 27)
    \item \textbf{C10: Advanced Diagnostic Testing (Find the Flaw):} Identifying vacuous logic errors, converse fallacies, and division-by-zero in cases. (Problems 28, 29, 30)
\end{itemize}

\end{document}